\documentclass[12pt,a4paper]{article}

% =========================
% Thai XeLaTeX Setup
% =========================
\usepackage{fontspec}
\usepackage{polyglossia}
\setdefaultlanguage{thai}
\setotherlanguage{english}
\setmainfont{TH Sarabun New}
\newfontfamily\thaifont{TH Sarabun New}
\newfontfamily\englishfont{Times New Roman}

% =========================
% Packages
% =========================
\usepackage[a4paper,margin=1in]{geometry}
\usepackage{amsmath,amssymb,amsfonts}
\usepackage{array}
\usepackage{booktabs}
\usepackage{longtable}
\usepackage{multirow}
\usepackage{graphicx}
\usepackage{xcolor}
\usepackage{hyperref}
\usepackage{enumitem}
\usepackage{titlesec}
\usepackage{fancyhdr}
\usepackage{setspace}
\usepackage{caption}
\usepackage{float}
\usepackage{tikz}
\usetikzlibrary{arrows.meta,positioning,shapes.geometric,fit}

% =========================
% Formatting
% =========================
\setlength{\parindent}{0pt}
\setlength{\parskip}{0.5em}
\onehalfspacing
\hypersetup{
	colorlinks=true,
	linkcolor=blue!50!black,
	urlcolor=blue!60!black,
	citecolor=blue!50!black,
	pdftitle={รายงานสรุปเชิงลึกการทดสอบความปลอดภัย FrodoKEM v2},
	pdfauthor={OpenAI}
}

\titleformat{\section}{\large\bfseries}{\thesection.}{0.5em}{}
\titleformat{\subsection}{\normalsize\bfseries}{\thesubsection}{0.5em}{}
\titleformat{\subsubsection}{\normalsize\bfseries}{\thesubsubsection}{0.5em}{}

\pagestyle{fancy}
\fancyhf{}
\lhead{รายงานสรุปเชิงลึกการทดสอบความปลอดภัย FrodoKEM v2}
\rhead{\thepage}

% =========================
% Custom commands
% =========================
\newcommand{\term}[1]{\textbf{#1}}
\newcommand{\eng}[1]{\textenglish{#1}}
\newcommand{\important}[1]{\textcolor{red!60!black}{\textbf{#1}}}
\newcommand{\vect}[1]{\mathbf{#1}}
\newcommand{\mat}[1]{\mathbf{#1}}

\begin{document}
	
	\begin{titlepage}
		\centering
		{\Large รายงานสรุปเชิงลึกเพื่อการอ้างอิงและจัดทำสื่อประกอบการนำเสนอ\par}
		\vspace{1.2cm}
		{\huge\bfseries การประเมินความปลอดภัยของ \eng{FrodoKEM v2} ภายใต้กลไก \eng{Masked-Salting}\par}
		\vspace{0.8cm}
		{\LARGE รายงานเชิงวิเคราะห์จากแผนการทดสอบ 3 ระดับ\par}
		\vspace{1.3cm}
		{\large ฉบับภาษาไทยในรูปแบบ \eng{LaTeX} สำหรับใช้สร้าง \eng{PDF} และต่อยอดเป็น \eng{Presentation}\par}
		\vfill
		\begin{tabular}{>{\raggedright\arraybackslash}p{0.28\textwidth} p{0.62\textwidth}}
			หัวข้อหลัก & การวิเคราะห์ความซับซ้อน, \eng{Exhaustive Search}, ช่องว่างประสิทธิภาพเชิงอัลกอริทึม, และการโจมตีแบบผสมผสาน \\
			กรอบระบบ & \eng{FrodoKEM v2} พร้อมกลไก \eng{Masked-Salting} \\
			วัตถุประสงค์ & สรุปรวบยอดข้อมูลทั้งหมดให้อยู่ในรูปเอกสารอ้างอิงที่ครบถ้วนสำหรับการเขียนรายงานหรือทำสไลด์ \\
			เอกสารฉบับนี้ & จัดโครงสร้างใหม่ให้เป็นรายงานเชิงเทคนิคฉบับสมบูรณ์ \\
		\end{tabular}
		\vfill
		{\large วันที่จัดทำ: \today\par}
	\end{titlepage}
	
	\tableofcontents
	\newpage
	
	% =========================================================
	\section{บทสรุปผู้บริหาร}
	
	รายงานฉบับนี้สังเคราะห์สาระจากเอกสารและอินโฟกราฟิกทั้งสี่ชุดที่มุ่งประเมินความมั่นคงของ \eng{FrodoKEM v2} หลังมีการเพิ่มกลไก \eng{Masked-Salting} เข้าไปในโครงสร้างเมทริกซ์หลักของระบบ โดยแก่นของการเปลี่ยนแปลงคือการเปลี่ยนจากเมทริกซ์ฐานเดิม $A$ ไปเป็นเมทริกซ์ที่ถูกปรับแล้วในรูป
	\[
	A' = A + M \pmod q
	\]
	เมื่อ $M$ เป็นเมทริกซ์แบบเบาบาง \eng{(sparse matrix)} ที่ฉีดค่าเพียง \term{1 จุดต่อ 1 คอลัมน์} เพื่อสร้างความไม่แน่นอนเพิ่มเติมในโครงสร้างปัญหาเชิงแลตทิซ
	
	ภาพรวมของข้อมูลทั้งหมดชี้ไปในทิศทางเดียวกันว่า การปรับปรุงนี้ไม่ได้ทำให้ระบบเกิดจุดอ่อนเชิงโครงสร้างที่ชัดเจน และไม่ได้สร้าง ``ทางลัด'' ให้กับผู้โจมตี ทั้งในกรณีของการเดาแบบครอบจักรวาล \eng{(exhaustive / brute-force)}, การลดรูปแลตทิซด้วย \eng{BKZ}, และการโจมตีแบบผสมผสานที่ใช้ \eng{BKZ} ช่วยตัดขอบเขตการค้นหาก่อนจึงตามด้วยการเดาแบบสุ่ม
	
	สาระสำคัญจากการวิเคราะห์สรุปได้ดังนี้
	\begin{enumerate}[label=\arabic*.]
		\item การทดสอบถูกออกแบบเป็น \term{สามระยะ} เพื่อให้ค่อย ๆ เพิ่มความเข้มข้นของแบบจำลองการโจมตี ตั้งแต่การพิสูจน์ความซับซ้อนพื้นฐาน ไปจนถึงการประเมินความได้เปรียบจากอัลกอริทึมขั้นสูง
		\item ระดับการทดลองใช้แนวทาง \term{Scale-down Testing} เพราะการทดสอบตรงกับพารามิเตอร์จริงระดับ $640\times640$ ยังเกินขอบเขตการคำนวณที่ทำได้จริงในงานประเมินเบื้องต้น
		\item ในทุกแผนการทดสอบ แนวโน้มหลักที่คาดหวังคือความซับซ้อนยังเติบโตแบบ \term{exponential} แม้จะพยายามใช้เทคนิคขั้นสูงเข้ามาช่วยลดขอบเขตการค้นหา
		\item เอกสารแต่ละชุดมีบทบาทเสริมกันอย่างชัดเจน: เอกสารแรกกำหนดกรอบของ \eng{Exhaustive Search}, เอกสารที่สองเปรียบเทียบช่องว่างระหว่าง \eng{Brute-force} กับ \eng{BKZ}, เอกสารที่สามประเมินการโจมตีแบบ \eng{Hybrid Attack}, และเอกสาร \eng{PDF} ทำหน้าที่สรุปภาพรวมเชิงวิธีวิทยาและสมมติฐาน
		\item หากแนวโน้มจากการทดลองระดับเล็กยังคงรูปเมื่อขยายไปยังมิติใช้งานจริง การโจมตีจะมีต้นทุนสูงระดับที่ไม่สามารถทำได้จริงในทางปฏิบัติ
	\end{enumerate}
	
	ดังนั้น รายงานนี้จึงสรุปว่า \eng{FrodoKEM v2} ภายใต้กลไก \eng{Masked-Salting} ยังรักษาขอบเขตความปลอดภัยไว้ได้อย่างแข็งแรงตามกรอบสมมติฐานที่กำหนด โดยมีสาระเพียงพอสำหรับใช้เป็นต้นฉบับอ้างอิงเพื่อจัดทำรายงานวิจัย บทความภายในองค์กร หรือสร้าง \eng{presentation} เชิงเทคนิคได้อย่างเป็นระบบ
	
	% =========================================================
	\section{ภูมิหลังและบริบทของปัญหา}
	
	ระบบเข้ารหัสหลังยุคควอนตัมที่อาศัยปัญหาเชิงแลตทิซ เช่น \eng{Learning With Errors (LWE)} หรือปัญหาที่เกี่ยวข้อง มีจุดแข็งอยู่ที่การทำให้การกู้คืนความลับจากข้อมูลสาธารณะเป็นปัญหาที่มีความยากสูงมากเชิงคำนวณ อย่างไรก็ตาม เมื่อมีการปรับแต่งโครงสร้างภายในของระบบ เช่น การเพิ่มกลไก \eng{Masked-Salting} ลงไปในเมทริกซ์สาธารณะ จำเป็นต้องตรวจสอบอย่างเป็นระบบว่าการเปลี่ยนแปลงดังกล่าว
	\begin{itemize}
		\item ไม่สร้าง \term{รูปแบบพิเศษ} ที่ผู้โจมตีใช้เป็นเบาะแสได้ง่ายขึ้น
		\item ไม่ทำให้เกิด \term{แลตทิซอ่อน} \eng{(weak lattices)} ซึ่งลดต้นทุนของ \eng{BKZ} ลงอย่างมีนัยสำคัญ
		\item ไม่เปิดช่องให้การโจมตีแบบผสมผสานตัดขนาด \eng{search space} ลงจนเหลือน้อยเกินไป
	\end{itemize}
	
	แนวคิดของ \eng{Masked-Salting} ในรายงานนี้คือ การเพิ่มเมทริกซ์ $M$ ที่มีความเบาบางสูงเข้าไปใน $A$ เพื่อสร้างเมทริกซ์ใหม่ $A'$ โดยยังคงทำงานภายใต้โมดูลัส $q$ จุดสำคัญคือการเลือกให้ $M$ มีโครงสร้าง ``กระจาย'' เพียง 1 จุดต่อ 1 คอลัมน์ ซึ่งตามสมมติฐานแล้วจะเพิ่มความไม่แน่นอน โดยไม่ทำให้ระบบมีอคติทางสถิติที่สังเกตได้ง่าย
	
	การตรวจสอบจึงไม่ได้มุ่งเพียงคำถามว่า ``เดายากหรือไม่'' แต่รวมถึงคำถามลึกกว่านั้น เช่น
	\begin{enumerate}[label=\alph*)]
		\item ผู้โจมตีสามารถแยกความแตกต่างระหว่างระบบรุ่นเดิมกับรุ่นที่มี \eng{Masked-Salting} ได้หรือไม่
		\item อัลกอริทึมลดรูปแลตทิซ เช่น \eng{BKZ} ได้รับประโยชน์จากโครงสร้างใหม่นี้หรือไม่
		\item หากใช้ \eng{BKZ} เป็นตัวช่วยคัดกรองก่อนแล้วค่อยเดาต่อ ความยากจะลดลงมากพอให้โจมตีได้จริงหรือไม่
	\end{enumerate}
	
	คำถามทั้งสามข้อข้างต้นกลายเป็นที่มาของกรอบวิธีวิทยาแบบสามระยะ ซึ่งเป็นแกนกลางของรายงานทั้งฉบับ
	
	% =========================================================
	\section{วัตถุประสงค์ของการประเมิน}
	
	การรวบรวมข้อมูลทั้งสี่แหล่งสะท้อนวัตถุประสงค์ร่วมที่ชัดเจนดังต่อไปนี้
	\begin{enumerate}[label=\arabic*.]
		\item \term{พิสูจน์ความทนทานของระบบต่อการเดาแบบครอบจักรวาล} ว่าการเพิ่ม \eng{Masked-Salting} ไม่ได้ลดระดับความยากพื้นฐานของปัญหา
		\item \term{วิเคราะห์ช่องว่างระหว่างความยากทางทฤษฎีกับการโจมตีจริง} โดยเปรียบเทียบระหว่าง \eng{Exhaustive Search} กับ \eng{BKZ Lattice Reduction}
		\item \term{ทดสอบสถานการณ์โจมตีที่เข้มข้นขึ้น} ด้วยการใช้ \eng{BKZ} ช่วยบีบขอบเขต แล้วจึงใช้ \eng{Brute-force} เฉพาะในส่วนที่เหลือ
		\item \term{ตรวจสอบพฤติกรรมเชิงแนวโน้ม} ว่าเมื่อมิติเพิ่มขึ้น ความยากยังเติบโตแบบเอ็กซ์โปเนนเชียล และไม่มีการลดรูปไปเป็นเชิงพหุนาม
		\item \term{สร้างหลักฐานเชิงอ้างอิง} สำหรับใช้สื่อสารในรูปแบบรายงานหรือสไลด์ โดยสรุปทั้งตรรกะการออกแบบ วิธีทดสอบ และสาระที่ตีความได้จากผลการทดลอง
	\end{enumerate}
	
	กล่าวอีกนัยหนึ่ง รายงานนี้ไม่ได้เสนอ ``ผลทดลองเชิงตัวเลขสุดท้าย'' ของพารามิเตอร์จริง แต่เสนอ \term{กรอบพิสูจน์ความแข็งแรงเชิงระบบ} ที่ใช้ข้อมูลระดับย่อส่วนเพื่อรองรับข้อสรุปเชิงแนวโน้ม
	
	% =========================================================
	\section{โครงสร้างการทดสอบแบบสามระยะ}
	
	การประเมินทั้งหมดถูกแบ่งเป็นสามระยะหลัก ดังภาพรวมในรูปแบบเชิงแนวคิดต่อไปนี้
	
	\begin{center}
		\begin{tikzpicture}[
			node distance=1.8cm and 1.4cm,
			box/.style={draw, rounded corners, thick, align=center, minimum width=0.26\textwidth, minimum height=1.8cm, fill=blue!5},
			arrow/.style={-{Latex[length=3mm]}, thick}
			]
			\node[box] (p1) {ระยะที่ 1\\\textbf{Exhaustive Search}\\วัดความซับซ้อนพื้นฐานและตรวจสอบ \eng{Indistinguishability}};
			\node[box, right=of p1] (p2) {ระยะที่ 2\\\textbf{Algorithm Efficiency Gap}\\เปรียบเทียบ \eng{Brute-force} กับ \eng{BKZ}};
			\node[box, right=of p2] (p3) {ระยะที่ 3\\\textbf{Hybrid Attack}\\ใช้ \eng{BKZ} ช่วย \eng{Pruning} ก่อน \eng{Brute-force}};
			\draw[arrow] (p1) -- (p2);
			\draw[arrow] (p2) -- (p3);
		\end{tikzpicture}
	\end{center}
	
	ทั้งสามระยะมีลักษณะต่อเนื่องกันเชิงตรรกะ กล่าวคือ
	\begin{itemize}
		\item ระยะที่ 1 ใช้พิสูจน์ว่าในระดับพื้นฐานแล้วระบบไม่ให้ ``ทางลัด'' จากการเดาแบบทื่อ ๆ
		\item ระยะที่ 2 ถามต่อว่า แม้มีอัลกอริทึมคณิตศาสตร์ขั้นสูง เช่น \eng{BKZ} เข้ามา ระบบยังคงปลอดภัยเพียงใด
		\item ระยะที่ 3 เป็นสถานการณ์เข้มที่สุดในเอกสารชุดนี้ โดยถือว่าผู้โจมตีฉลาดพอที่จะใช้ \eng{BKZ} ช่วยลดงานบางส่วนก่อน แล้วค่อยลงมือเดาในส่วนที่เหลือ
	\end{itemize}
	
	หากทั้งสามระยะยังให้ภาพสอดคล้องกันว่า ``ความซับซ้อนยังสูงมาก'' ก็จะกลายเป็นหลักฐานเชิงระบบที่หนักแน่นกว่าการดูเฉพาะการทดลองแบบใดแบบหนึ่งเพียงลำพัง
	
	% =========================================================
	\section{กลยุทธ์การลดขนาดปัญหา (Scale-down Testing)}
	
	เอกสารทั้งหมดเห็นตรงกันว่าการโจมตีพารามิเตอร์จริงของระบบทันที เช่น มิติระดับ $640\times640$ ไม่เหมาะกับการทดสอบเบื้องต้น เพราะมีต้นทุนด้านเวลาและทรัพยากรสูงมาก จึงใช้แนวทาง \term{Scale-down Testing} แทน โดยแบ่งระดับดังนี้
	
	\begin{table}[H]
		\centering
		\caption{ระดับการทดสอบที่ใช้ในรายงาน}
		\begin{tabular}{p{0.18\textwidth} p{0.23\textwidth} p{0.49\textwidth}}
			\toprule
			ระดับ & ช่วงมิติ/พารามิเตอร์ & วัตถุประสงค์หลัก \\
			\midrule
			\eng{Nano Scale} & $2\times2$ และค่าพารามิเตอร์เล็ก & ตรวจสอบความถูกต้องของ \eng{logic}, \eng{data path}, การสร้างเมทริกซ์ และการตรวจสอบเชิงโครงสร้าง \\
			\eng{Micro Scale} & โดยประมาณ $3\times3$ ถึง $6\times6$ & เก็บจุดข้อมูลเพื่อดูแนวโน้มของเวลา ความซับซ้อน และประสิทธิภาพของอัลกอริทึมเมื่อมิติเริ่มขยาย \\
			\eng{Mini Scale} & ประมาณ $8\times8$ & ทดสอบขีดจำกัดระดับเล็กสุดที่ยังพอประมวลผลได้จริงในเวลาเหมาะสม และใช้เป็นจุดเริ่มต้นของการประมาณแนวโน้ม \\
			\bottomrule
		\end{tabular}
	\end{table}
	
	เหตุผลสำคัญของแนวทางนี้มีสองประการ
	\begin{enumerate}[label=\arabic*.]
		\item ต้องการเห็น \term{รูปแบบการเติบโต} ของความซับซ้อน ไม่ใช่เพียงค่าตัวเลขเฉพาะจุด
		\item ต้องการแยกแยะว่า ``ความยาก'' ที่สังเกตได้เกิดจากโครงสร้างคณิตศาสตร์ของระบบจริง หรือเป็นเพียงข้อจำกัดจากทรัพยากรทดลอง
	\end{enumerate}
	
	ด้วยเหตุนี้ การทดลองระดับเล็กจึงไม่ได้มีบทบาทเป็นเพียง ``การทดสอบย่อ'' แต่ทำหน้าที่เป็น \term{กล้องขยายเชิงพฤติกรรม} เพื่อดูว่าระบบโตขึ้นอย่างไรเมื่อมิติขยายขึ้นทีละขั้น
	
	% =========================================================
	\section{ระยะที่ 1: การวิเคราะห์แบบ Exhaustive Search}
	
	\subsection{แนวคิดหลัก}
	
	ระยะที่ 1 ใช้การเดาแบบครอบจักรวาลหรือ \eng{Exhaustive Search} เป็นฐานอ้างอิงของความซับซ้อน โดยวัตถุประสงค์สำคัญไม่ใช่เพียงการถามว่า ``เดาได้หรือไม่'' แต่คือการตรวจสอบว่า \eng{Masked-Salting} ทำให้การกระจายของระบบผิดไปจากเดิมหรือไม่
	
	หากเมทริกซ์ใหม่ $A'$ ยังคงมีพฤติกรรมเชิงสถิติคล้ายเมทริกซ์เดิม $A$ การเดาเพื่อแยกองค์ประกอบของ $M$ ออกจากระบบย่อมไม่ง่ายขึ้นอย่างมีนัยสำคัญ ซึ่งสอดคล้องกับแนวคิด \eng{Indistinguishability}
	
	\subsection{สูตรความซับซ้อนพื้นฐาน}
	
	จากข้อมูลในเอกสาร แนวคิดของพื้นที่ค้นหา \eng{search space} สำหรับการเดาแบบพื้นฐานถูกสรุปในรูป
	\[
	W = (n \times q)^n
	\]
	เมื่อ $n$ คือมิติของปัญหา และ $q$ คือโมดูลัสที่เกี่ยวข้องกับโครงสร้างของเมทริกซ์/ข้อผิดพลาด
	
	สูตรนี้ทำหน้าที่เป็นกรอบเชิงสัญลักษณ์เพื่อบอกว่า
	\begin{itemize}
		\item เมื่อ $n$ เพิ่มขึ้น ความยากไม่ได้เพิ่มแบบเชิงเส้น
		\item เมื่อ $q$ เพิ่มขึ้น จำนวนทางเลือกต่อมิติก็เพิ่มขึ้น
		\item เมื่อทั้งสองค่าขยายพร้อมกัน ระบบจะเติบโตเร็วแบบ \term{ทวีคูณ} หรือ \term{exponential}
	\end{itemize}
	
	\subsection{สาระสำคัญของการทดสอบ}
	
	เอกสารเกี่ยวกับ \eng{Exhaustive Search} เน้นประเด็นสำคัญดังนี้
	\begin{enumerate}[label=\arabic*.]
		\item ใช้การทดสอบระดับเล็กเป็นตัวแทนเพื่อตรวจสอบเส้นแนวโน้มของความซับซ้อน
		\item เปรียบเทียบพฤติกรรมระหว่างระบบต้นฉบับกับระบบที่มี \eng{Masked-Salting}
		\item พิจารณาว่าเวลาเฉลี่ยในการเดาและรูปแบบการเติบโตยังคงคล้ายกันหรือไม่
		\item ใช้ข้อสังเกตนี้เป็นฐานสำหรับอ้างว่า \eng{Masked-Salting} ไม่ได้เพิ่มอคติที่ทำให้ผู้โจมตีได้ประโยชน์
	\end{enumerate}
	
	\subsection{ขั้นตอนเชิงตรรกะของการทดลอง}
	
	ในเชิงแนวคิด สามารถสรุปกระบวนการของระยะที่ 1 ได้เป็นลำดับดังนี้
	\begin{enumerate}[label=\arabic*.]
		\item สร้างตัวอย่างข้อมูลในระดับพารามิเตอร์เล็ก ทั้งกรณีไม่มีและมี \eng{Masked-Salting}
		\item ให้ผู้โจมตีสมมุติพยายามเดาความแตกต่างหรือองค์ประกอบของ $M$
		\item วัดจำนวนครั้งที่ต้องลองหรือเวลาที่ใช้ก่อนการตรวจสอบจะผ่าน
		\item เปรียบเทียบผลของเวอร์ชันเดิมกับเวอร์ชันที่ถูกปรับโครงสร้าง
		\item สรุปว่าเวลาเฉลี่ยและอัตราการเติบโตบ่งชี้การแยกแยะได้หรือไม่
	\end{enumerate}
	
	\subsection{ข้อสรุปเชิงตีความของระยะที่ 1}
	
	ถ้าระบบ \eng{v2} แสดงรูปแบบเวลาและแนวโน้มการเติบโตใกล้เคียงกับระบบเดิม \eng{v1} อย่างต่อเนื่อง จะตีความได้ว่า
	\begin{itemize}
		\item การเพิ่ม $M$ ไม่ได้สร้างลายเซ็นเชิงสถิติที่ผู้โจมตีจับได้ง่าย
		\item พื้นที่ค้นหายังคงใหญ่ในอัตราที่สอดคล้องกับปัญหาที่ยากอยู่เดิม
		\item การโจมตีแบบเดาทื่อ ๆ ไม่มี ``ทางลัด'' ที่เห็นได้จากโครงสร้างที่เพิ่มเข้ามา
	\end{itemize}
	
	กล่าวอย่างย่อ ระยะที่ 1 ทำหน้าที่ตอบคำถามพื้นฐานที่สุดว่า ``แม้โจมตีแบบง่ายที่สุด ระบบยังคงดูเหมือนเดิมหรือไม่''
	
	% =========================================================
	\section{ระยะที่ 2: การวิเคราะห์ช่องว่างประสิทธิภาพเชิงอัลกอริทึม}
	
	\subsection{เหตุผลที่ต้องมีระยะที่ 2}
	
	การพิสูจน์ว่าระบบปลอดภัยต่อ \eng{Brute-force} เพียงอย่างเดียวอาจยังไม่เพียงพอ เพราะในโลกจริง ผู้โจมตีจะไม่เลือกวิธีที่โง่ที่สุดเสมอไป หากมีอัลกอริทึมเชิงคณิตศาสตร์ที่ช่วยลดภาระการค้นหาได้ ก็จำเป็นต้องประเมินด้วยว่าระบบต้านทานได้ดีเพียงใด
	
	จึงเกิดระยะที่ 2 ขึ้นเพื่อวิเคราะห์ \term{Algorithm Efficiency Gap} หรือ ``ช่องว่างระหว่างความยากตามทฤษฎีกับความเร็วจริงของอัลกอริทึมขั้นสูง'' โดยเน้นที่การใช้ \eng{BKZ Lattice Reduction}
	
	\subsection{แนวคิดของช่องว่างประสิทธิภาพ}
	
	เอกสารชุดนี้เสนอให้เปรียบเทียบระหว่าง
	\begin{itemize}
		\item เวลาหรือความซับซ้อนของ \eng{Exhaustive Search}
		\item เวลาหรือความซับซ้อนของ \eng{BKZ} หลังสร้างแลตทิซที่สะท้อนโครงสร้างของระบบ
	\end{itemize}
	
	โดยแนวคิดหลักคือ หาก \eng{BKZ} เร็วกว่าการเดาทั่วไปมาก ก็แปลว่าผู้โจมตีมี ``ความได้เปรียบเชิงอัลกอริทึม'' แต่คำถามสำคัญคือความได้เปรียบนั้นมากพอจะเปลี่ยนระบบจาก ``โจมตีไม่ได้'' ให้กลายเป็น ``โจมตีได้จริง'' หรือไม่
	
	\subsection{การปรับ \eng{BKZ} ให้เหมาะกับ \eng{Masked-Salting}}
	
	จุดเด่นของเอกสารนี้คือไม่ได้มอง \eng{BKZ} แบบสำเร็จรูป แต่พูดถึงความจำเป็นในการปรับฐานแลตทิซให้รองรับโครงสร้าง \eng{Sparse Salting} โดยมีองค์ประกอบแนวคิดสำคัญดังนี้
	\begin{enumerate}[label=\arabic*.]
		\item สร้าง \term{embedding} ที่เหมาะกับเมทริกซ์ $A'$ แทนที่จะใช้เฉพาะโครงสร้างมาตรฐานของปัญหาเดิม
		\item พิจารณาการเลือกค่า \eng{block size} หรือ $\beta$ ของ \eng{BKZ} เพื่อดูว่าที่ระดับมิติต่าง ๆ ผู้โจมตีได้ประโยชน์มากน้อยเพียงใด
		\item คำนึงถึงการทำงานภายใต้โมดูลัส $q$ และโครงสร้างของ \eng{$q$-ary lattices}
	\end{enumerate}
	
	ในเชิงแนวคิด รายงานต้นทางยกตัวอย่างการสร้างฐานในลักษณะคล้าย
	\[
	\mat{B} = \begin{bmatrix} I_n & A'^T \\ 0 & qI_n \end{bmatrix}
	\]
	ก่อนจะนำฐานดังกล่าวไปลดรูปด้วย \eng{BKZ} เพื่อสืบค้นเวกเตอร์สั้นที่อาจสัมพันธ์กับความลับหรือโครงสร้างบางส่วนของระบบ
	
	\subsection{ตัวชี้วัดที่ใช้ตีความ}
	
	สาระที่ถูกเน้นในเอกสารมีอยู่สองส่วนสำคัญ
	\begin{enumerate}[label=\arabic*.]
		\item \term{อัตราส่วนประสิทธิภาพ} ระหว่าง \eng{Exhaustive Search} กับ \eng{BKZ} เพื่อดูว่าผู้โจมตีได้เปรียบเพียงใด
		\item \term{Root Hermite Factor} และความแข็งแรงของฐานแลตทิซ เพื่อประเมินว่าโครงสร้างที่เติมเข้าไปสร้างจุดอ่อนหรือไม่
	\end{enumerate}
	
	ถ้า \eng{BKZ} แม้จะเร็วขึ้นกว่าการเดาทื่อ ๆ แต่ยังเติบโตแบบเอ็กซ์โปเนนเชียล และยังไม่ทำให้ปัญหาลดลงเป็นงานที่ทำได้จริง ก็จะสรุปได้ว่าความได้เปรียบดังกล่าวยังอยู่ในระดับจำกัด
	
	\subsection{แก่นของข้อสรุปจากระยะที่ 2}
	
	ข้อมูลในเอกสารชี้ไปในทิศทางว่า
	\begin{itemize}
		\item \eng{BKZ} อาจช่วยให้โจมตีได้ดีกว่าการเดาสุ่มในมิติเล็ก
		\item แต่ความได้เปรียบดังกล่าว \term{ไม่ได้เปลี่ยนธรรมชาติของความยาก} ให้กลายเป็นปัญหาเชิงพหุนาม
		\item เมื่อมิติเพิ่มขึ้น ช่องว่างระหว่าง ``โจมตีได้ในทางทฤษฎี'' กับ ``โจมตีได้จริงในทางปฏิบัติ'' ยังคงกว้างมาก
	\end{itemize}
	
	ด้วยเหตุนี้ ระยะที่ 2 จึงทำหน้าที่ตอบคำถามว่า ``ถึงแม้ผู้โจมตีจะฉลาดขึ้นและใช้คณิตศาสตร์เก่งขึ้น ระบบยังแข็งแรงอยู่หรือไม่''
	
	% =========================================================
	\section{ระยะที่ 3: การโจมตีแบบ Hybrid Attack}
	
	\subsection{แนวคิดของการโจมตีแบบผสมผสาน}
	
	ระยะที่ 3 ถือเป็นสถานการณ์โจมตีที่เข้มที่สุดในข้อมูลชุดนี้ เพราะไม่สมมุติว่าผู้โจมตีจะพึ่ง \eng{Brute-force} ล้วน หรือ \eng{BKZ} ล้วน แต่ให้ทั้งสองวิธีทำงานร่วมกัน กล่าวคือ
	\begin{enumerate}[label=\arabic*.]
		\item ใช้ \eng{BKZ} เพื่อช่วยคัดกรองหรือ \eng{prune} พิกัดที่มีแนวโน้มเป็นไปได้มากกว่า
		\item จากนั้นใช้ \eng{Exhaustive Search} เฉพาะในขอบเขตที่เหลือซึ่งถูกบีบให้เล็กลงแล้ว
	\end{enumerate}
	
	รูปแบบนี้เป็นสิ่งที่ควรกลัวมากกว่าการเดาล้วน ๆ เพราะถ้า \eng{BKZ} ช่วยตัดพื้นที่ค้นหาได้มากพอ ความซับซ้อนที่เหลืออาจลดลงจนโจมตีได้จริง
	
	\subsection{กระบวนการของ \eng{Hybrid Attack}}
	
	เอกสารต้นทางสรุปแนวคิดของกระบวนการนี้เป็น 3 ช่วงดังนี้
	\begin{enumerate}[label=\arabic*.]
		\item \term{Lattice Embedding}: สร้างแบบจำลองแลตทิซจาก $A'$ และพารามิเตอร์ที่เกี่ยวข้อง
		\item \term{BKZ Partial Solve}: ใช้ \eng{BKZ} เพื่อลดรูปฐานและคัดกรองพิกัดหรือโครงสร้างที่ดูมีความเป็นไปได้สูง
		\item \term{Final Brute-force}: เดาเฉพาะค่าที่เหลือในขอบเขตที่ถูกลดลงแล้ว
	\end{enumerate}
	
	\subsection{คำถามหลักของระยะนี้}
	
	ระยะที่ 3 ไม่ได้ถามว่า \eng{BKZ} ดีหรือไม่ดี แต่ถามลึกลงไปอีกว่า
	\begin{quote}
		``หาก \eng{BKZ} ช่วยตัดงานบางส่วนออกไปแล้ว พื้นที่ค้นหาที่เหลือยังใหญ่พอจะปกป้องระบบอยู่หรือไม่''
	\end{quote}
	
	นี่คือแก่นของการประเมิน \eng{Hybrid Attack}
	
	\subsection{สาระเชิงตีความจากข้อมูลต้นทาง}
	
	ข้อมูลชุดนี้สื่อข้อสรุปสำคัญ 3 ประการ
	\begin{enumerate}[label=\arabic*.]
		\item ในมิติเล็ก \eng{BKZ} อาจช่วยลด \eng{search space} ได้บ้าง
		\item แต่การลดนั้นยังมีลักษณะเป็นเพียง \term{การลดค่าคงที่หรือเปอร์เซ็นต์บางส่วน} ไม่ใช่การเปลี่ยนรูปแบบการเติบโตของความซับซ้อน
		\item โครงสร้าง \eng{1 จุดต่อ 1 คอลัมน์} ของ $M$ มีความกระจายและไม่สร้างความสัมพันธ์เชิงพื้นที่ที่ง่ายต่อการทำนาย ทำให้ \eng{BKZ} ไม่สามารถคาดเดาล่วงหน้าในคอลัมน์ถัดไปได้ดีนัก
	\end{enumerate}
	
	กล่าวอย่างตรงประเด็นคือ \eng{Hybrid Attack} ทำให้ผู้โจมตี ``ฉลาดขึ้น'' แต่ยังไม่ทำให้การโจมตี ``เป็นไปได้จริง''
	
	% =========================================================
	\section{การเปรียบเทียบสาระจากทั้งสี่เอกสาร}
	
	เมื่อพิจารณาเอกสารทั้งสี่ชุดร่วมกัน จะพบว่ามีความสัมพันธ์เชิงบทบาทดังตารางต่อไปนี้
	
	\begin{longtable}{p{0.21\textwidth} p{0.26\textwidth} p{0.44\textwidth}}
		\caption{บทบาทของเอกสารแต่ละชุดในภาพรวมการประเมิน}\\
		\toprule
		เอกสาร & ประเด็นหลัก & บทบาทต่อรายงานรวม \\
		\midrule
		\endfirsthead
		\toprule
		เอกสาร & ประเด็นหลัก & บทบาทต่อรายงานรวม \\
		\midrule
		\endhead
		\eng{Exhaustion\_Search\_Test.html} & การพิสูจน์ความซับซ้อนพื้นฐานและแนวคิด \eng{Scale-down} & วางฐานให้เห็นว่าการเดาแบบครอบจักรวาลยังเติบโตแบบ \eng{exponential} และใช้เป็น \eng{baseline} ของทั้งระบบ \\
		\eng{Algorithm\_Efficiency\_Gap.html} & เปรียบเทียบความยากของ \eng{Brute-force} กับ \eng{BKZ} & ทำหน้าที่ประเมินความได้เปรียบของผู้โจมตีเมื่อใช้อัลกอริทึมขั้นสูง และตรวจสอบว่ามี \eng{weak lattice} หรือไม่ \\
		\eng{Lattice-Reduction\_Aided\_Brute-force.html} & การโจมตีแบบผสมผสาน \eng{BKZ + Brute-force} & เป็นกรณีเข้มที่สุด โดยวิเคราะห์ว่าหลังการ \eng{pruning} แล้วความซับซ้อนที่เหลือยังมากเพียงใด \\
		\eng{PDF Summary Report} & บทสรุปวิธีวิทยาและสมมติฐานรวม & ทำหน้าที่เชื่อมทั้งสามระยะเข้าด้วยกัน และสรุปสาระในภาพรวมเดียว \\
		\bottomrule
	\end{longtable}
	
	ข้อดีของการมีหลายเอกสารคือทำให้การสรุปครั้งนี้ไม่พึ่งข้อมูลเชิงมุมมองเดียว แต่ค่อย ๆ ต่อภาพจาก
	\begin{itemize}
		\item การตั้งคำถามพื้นฐาน,
		\item ไปสู่การเปรียบเทียบเชิงอัลกอริทึม,
		\item และจบที่สถานการณ์โจมตีแบบผสมที่รุนแรงกว่า
	\end{itemize}
	
	นี่ทำให้รายงานรวมมีน้ำหนักสำหรับใช้เป็นต้นทางอ้างอิงในการอธิบายต่อผู้อ่านหรือผู้ฟังในงานนำเสนอ
	
	% =========================================================
	\section{สมมติฐานสำคัญที่ใช้ตลอดการวิเคราะห์}
	
	จากการสังเคราะห์ข้อมูลทั้งหมด สมมติฐานหลักสามารถเรียบเรียงใหม่ให้ชัดเจนได้ดังนี้
	
	\subsection{สมมติฐานที่ 1: ความซับซ้อนยังเติบโตแบบ Exponential}
	
	เมื่อมิติ $n$ เพิ่มขึ้น เวลาและต้นทุนในการโจมตีจะเพิ่มขึ้นอย่างรวดเร็วในลักษณะที่เมื่อแสดงบนกราฟ \eng{log-scale} แล้วควรมีแนวโน้มใกล้เส้นตรง ซึ่งสื่อถึงการเติบโตแบบเอ็กซ์โปเนนเชียลบนสเกลปกติ
	
	\subsection{สมมติฐานที่ 2: ระบบยังคง Indistinguishable}
	
	เมทริกซ์ $A'$ ที่ได้จากการเพิ่ม $M$ จะยังคงไม่สามารถแยกจากกรณีเดิมได้ง่ายในเชิงสถิติ กล่าวคือผู้โจมตีไม่ควรได้เบาะแสจาก ``รูปร่าง'' การกระจายของข้อมูลที่เปลี่ยนไป
	
	\subsection{สมมติฐานที่ 3: ไม่มี Weak Lattice}
	
	การปรับโครงสร้างของเมทริกซ์ไม่ควรสร้างฐานแลตทิซที่ลดรูปได้ง่ายเป็นพิเศษ หากเกิด \eng{weak lattice} ขึ้นจริง \eng{BKZ} อาจได้เปรียบเกินคาดและลดต้นทุนของการโจมตีลงอย่างมีนัยสำคัญ
	
	\subsection{สมมติฐานที่ 4: Hybrid Attack ยังไม่เพียงพอ}
	
	แม้จะมีการคัดกรองพิกัดล่วงหน้าด้วย \eng{BKZ} พื้นที่ค้นหาที่เหลือยังต้องโตแบบเอ็กซ์โปเนนเชียลอยู่ดี จึงไม่ควรนำไปสู่การโจมตีที่ทำได้จริงในระดับพารามิเตอร์ใช้งาน
	
	\subsection{สมมติฐานที่ 5: การขยายแนวโน้มไปยังระดับจริงยังคงให้ภาพ ``โจมตีไม่ได้ในทางปฏิบัติ''}
	
	แม้รายงานจะทดสอบในระดับมิติเล็ก แต่หากรูปแบบการเติบโตคงเส้นคงวา เมื่อขยายไปสู่ระดับจริง เช่น $640\times640$ ต้นทุนจะสูงเกินกว่าขีดจำกัดของระบบคอมพิวเตอร์ทั่วไปอย่างมาก
	
	% =========================================================
	\section{การตีความกราฟและภาพประกอบในเอกสารต้นทาง}
	
	เอกสารทั้งสามชุดในรูปแบบ \eng{HTML infographic} มีกราฟประกอบหลายชนิด เช่น \eng{line chart}, \eng{bar chart}, \eng{bubble chart} และแผนผังการไหลเชิงโครงสร้าง แม้ตัวเลขในกราฟจะทำหน้าที่เชิงสาธิตมากกว่าผลทดลองจริงในทุกกรณี แต่สามารถตีความเป็นสาระเพื่อใช้ทำสไลด์ได้ดังนี้
	
	\subsection{กราฟความซับซ้อนตามมิติ}
	
	กราฟชนิดนี้สื่อสารประเด็นสำคัญมากที่สุด คือเมื่อ $n$ เพิ่มขึ้น ความซับซ้อนเพิ่มเร็วมากจนการประเมินบนแกนเชิงเส้นแทบไม่เห็นรายละเอียด จึงต้องใช้ \eng{logarithmic scale}
	
	\subsection{กราฟเปรียบเทียบ \eng{Brute-force} กับ \eng{BKZ}}
	
	กราฟเหล่านี้ช่วยสื่อว่า \eng{BKZ} อาจทำงานได้เร็วกว่า แต่ ``เร็วกว่า'' ไม่ได้แปลว่า ``เจาะได้จริง'' เพราะยังต้องดูด้วยว่าอัตราการเติบโตของเวลาในภาพรวมลดลงเป็นแบบใด หากยังเป็นเอ็กซ์โปเนนเชียลอยู่ ระบบก็ยังคงมีขอบเขตความปลอดภัยสูง
	
	\subsection{กราฟการลดขนาดของ Search Space ใน \eng{Hybrid Attack}}
	
	แม้การใช้ \eng{BKZ} ช่วยลดพื้นที่ค้นหาได้บางส่วน แต่กราฟพยายามสื่อว่าพื้นที่ที่เหลือยังคงใหญ่พอจะทำให้การโจมตีไม่มีความเป็นไปได้เชิงปฏิบัติ โดยเฉพาะเมื่อมิติเริ่มขยับขึ้น
	
	\subsection{กราฟแนวโน้มการทับซ้อนระหว่าง v1 และ v2}
	
	เอกสาร \eng{Exhaustive Search} เน้นประเด็นนี้ชัดเจน คือถ้าพฤติกรรมของ \eng{v1} และ \eng{v2} ซ้อนกันหรือใกล้เคียงกันมาก จะตีความได้ว่า \eng{Masked-Salting} ไม่ได้เปิดเผยรูปแบบใหม่ให้โจมตีได้ง่ายขึ้น
	
	% =========================================================
	\section{ข้อค้นพบเชิงเนื้อหาที่เหมาะสำหรับใช้ทำ Presentation}
	
	จากการรวมข้อมูลทั้งหมด สามารถแปลงเป็น ``แกนเรื่อง'' สำหรับงานนำเสนอได้อย่างมีประสิทธิภาพ โดยเสนอให้จัดประเด็นสำคัญเป็น 7 ส่วนดังนี้
	
	\subsection{ส่วนที่ 1: ปัญหาวิจัย}
	
	เริ่มด้วยคำถามหลักว่า การเพิ่ม \eng{Masked-Salting} ใน \eng{FrodoKEM v2} ทำให้ระบบแข็งแรงขึ้นโดยไม่สร้างจุดอ่อนใหม่ได้จริงหรือไม่
	
	\subsection{ส่วนที่ 2: นิยามการเปลี่ยนแปลงของระบบ}
	
	อธิบายสมการ
	\[
	A' = A + M \pmod q
	\]
	พร้อมเน้นว่า $M$ เป็น \eng{sparse matrix} ที่ฉีดค่าเพียง 1 จุดต่อคอลัมน์ จุดนี้ควรใช้ภาพหรือแอนิเมชันประกอบเพื่อให้เข้าใจได้เร็ว
	
	\subsection{ส่วนที่ 3: ทำไมต้องใช้ Scale-down Testing}
	
	อธิบายว่าพารามิเตอร์จริงใหญ่เกินไปสำหรับการทดลองตรง จึงใช้การวิเคราะห์ระดับย่อส่วนเพื่อจับพฤติกรรมการเติบโตของความซับซ้อน
	
	\subsection{ส่วนที่ 4: ระยะที่ 1 --- Exhaustive Search}
	
	ชี้ให้เห็นว่าเป็นการตั้งฐานว่า หาก \eng{v2} ยังดูเหมือน \eng{v1} ในเชิงสถิติ ผู้โจมตีก็ไม่ควรได้ประโยชน์จากการสังเกตโครงสร้างใหม่
	
	\subsection{ส่วนที่ 5: ระยะที่ 2 --- Efficiency Gap}
	
	อธิบายว่าระบบไม่ได้ถูกทดสอบแค่กับผู้โจมตีแบบ ``สุ่มเดา'' แต่ยังทดสอบกับผู้โจมตีที่ใช้คณิตศาสตร์แลตทิซระดับสูงด้วย เพื่อดูว่าความได้เปรียบเชิงอัลกอริทึมมีมากแค่ไหน
	
	\subsection{ส่วนที่ 6: ระยะที่ 3 --- Hybrid Attack}
	
	นำเสนอว่าแม้ผู้โจมตีจะฉลาดขึ้นโดยใช้ \eng{BKZ} ช่วยคัดกรองก่อน ความซับซ้อนที่เหลือก็ยังโตเร็วเกินกว่าจะโจมตีได้จริง
	
	\subsection{ส่วนที่ 7: ข้อสรุปใหญ่}
	
	ปิดท้ายด้วยประโยคที่ทรงพลัง เช่น
	\begin{quote}
		กลไก \eng{Masked-Salting} ใน \eng{FrodoKEM v2} ไม่ได้มอบ ``ทางลัด'' ให้ผู้โจมตี ทั้งในการเดาแบบครอบจักรวาล การโจมตีด้วย \eng{BKZ} หรือการโจมตีแบบผสมผสาน
	\end{quote}
	
	% =========================================================
	\section{ข้อจำกัดของข้อมูลและสิ่งที่ควรระวังในการอ้างอิง}
	
	แม้ข้อมูลทั้งหมดจะมีประโยชน์มากสำหรับใช้จัดทำรายงานหรือสไลด์ แต่ก็ควรระบุข้อจำกัดอย่างตรงไปตรงมาด้วย
	\begin{enumerate}[label=\arabic*.]
		\item เอกสารเชิง \eng{HTML infographic} หลายส่วนเป็นการนำเสนอเชิงแนวคิดและแนวโน้ม มากกว่าจะเป็นผลทดลองภาคสนามเต็มรูปแบบ
		\item ตัวเลขบางชุดในกราฟมีลักษณะเป็นค่าตัวอย่างเชิงสาธิต เพื่ออธิบายทิศทาง ไม่ใช่ค่ารับรองเชิงมาตรฐานจากการทดลองขนาดใหญ่ทั้งหมด
		\item การขยายผลจากมิติเล็กไปยังพารามิเตอร์จริงต้องอาศัยสมมติฐานว่าพฤติกรรมการเติบโตคงรูปต่อเนื่อง ซึ่งควรกล่าวอย่างระมัดระวังในสไลด์หรือรายงานวิจัย
		\item หากจะใช้ในงานอ้างอิงระดับวิชาการสูง ควรเสริมด้วยผลทดลองเชิงตัวเลขที่ได้จากการรันจริงในสภาพแวดล้อมควบคุม
	\end{enumerate}
	
	อย่างไรก็ตาม สำหรับจุดประสงค์ของการจัดทำรายงานอธิบายแนวคิด การสรุปวิธีวิทยา และการสร้างงานนำเสนอ เอกสารชุดนี้มีเนื้อหาครบถ้วนเพียงพอและมีลำดับเรื่องที่ดีมาก
	
	% =========================================================
	\section{ข้อเสนอแนะสำหรับการเรียบเรียงเป็นสไลด์นำเสนอ}
	
	เพื่อให้เอกสารฉบับนี้นำไปแปลงเป็น \eng{presentation} ได้ง่าย รายงานขอเสนอรูปแบบสไลด์ 10 หน้า ดังนี้
	
	\begin{longtable}{p{0.13\textwidth} p{0.25\textwidth} p{0.48\textwidth}}
		\caption{ข้อเสนอรูปแบบสไลด์นำเสนอ}\\
		\toprule
		สไลด์ & หัวข้อ & เนื้อหาที่ควรใส่ \\
		\midrule
		\endfirsthead
		\toprule
		สไลด์ & หัวข้อ & เนื้อหาที่ควรใส่ \\
		\midrule
		\endhead
		1 & ชื่อเรื่อง & ชื่อโครงการ, ระบบ \eng{FrodoKEM v2}, คำว่า \eng{Masked-Salting}, และประโยคชี้วัตถุประสงค์ \\
		2 & ภูมิหลัง & ปัญหาของการประเมินความปลอดภัยหลังปรับโครงสร้างเมทริกซ์ \\
		3 & กลไกหลัก & แสดงสมการ $A' = A + M \pmod q$ พร้อมคำอธิบายว่า $M$ เป็น \eng{sparse matrix} \\
		4 & วิธีวิทยา & แนะนำแนวทาง \eng{Scale-down Testing} และระดับ \eng{Nano/Micro/Mini} \\
		5 & ระยะที่ 1 & อธิบาย \eng{Exhaustive Search} และแนวคิด \eng{Indistinguishability} \\
		6 & ระยะที่ 2 & อธิบาย \eng{Algorithm Efficiency Gap} และบทบาทของ \eng{BKZ} \\
		7 & ระยะที่ 3 & อธิบาย \eng{Hybrid Attack} และเหตุผลว่าทำไมจึงเป็นสถานการณ์เข้มที่สุด \\
		8 & กราฟแนวโน้ม & ใส่กราฟความซับซ้อนบนแกน \eng{log-scale} และกราฟเปรียบเทียบ \eng{v1/v2} \\
		9 & ข้อค้นพบ & สรุปว่าระบบไม่เกิด \eng{weak lattice}, ไม่มีทางลัด, และยังโตแบบ \eng{exponential} \\
		10 & บทสรุป & เน้นสาระว่าระบบยัง \eng{robust} และโจมตีไม่ได้ในทางปฏิบัติ \\
		\bottomrule
	\end{longtable}
	
	% =========================================================
	\section{สรุปสุดท้าย}
	
	เมื่อสังเคราะห์ข้อมูลจากทั้งเอกสาร \eng{HTML} สามชุดและรายงานสรุป \eng{PDF} อีกหนึ่งชุด จะเห็นโครงเรื่องที่ชัดเจนมากว่า การประเมินความปลอดภัยของ \eng{FrodoKEM v2} ภายใต้กลไก \eng{Masked-Salting} ถูกออกแบบอย่างเป็นลำดับจากง่ายไปยาก และพยายามตอบคำถามสำคัญครบทุกมิติ ได้แก่
	\begin{itemize}
		\item การเดาแบบพื้นฐานยังยากอยู่หรือไม่
		\item อัลกอริทึมลดรูปแลตทิซช่วยผู้โจมตีได้มากแค่ไหน
		\item ถ้าใช้ทั้งคณิตศาสตร์และการเดาแบบผสมกันแล้ว ระบบยังปลอดภัยอยู่หรือไม่
	\end{itemize}
	
	ข้อสรุปเชิงรวมของเอกสารทั้งหมดคือ \term{ใช่ ระบบยังปลอดภัยอยู่} ภายใต้กรอบสมมติฐานที่ใช้ในงานนี้ โดยเฉพาะอย่างยิ่งไม่มีหลักฐานจากการวิเคราะห์ที่บ่งชี้ว่ากลไก \eng{Masked-Salting} สร้างจุดอ่อนใหม่ที่ชัดเจน หรือเปลี่ยนธรรมชาติของปัญหาให้โจมตีได้ง่ายขึ้นอย่างมีนัยสำคัญ
	
	กล่าวในภาษาที่เหมาะกับการนำเสนอ คือ
	\begin{quote}
		\textbf{\eng{Masked-Salting} ใน \eng{FrodoKEM v2} เพิ่มความซับซ้อนเชิงโครงสร้างโดยไม่เปิดเผยทางลัดให้ผู้โจมตี ทั้งในการโจมตีแบบเดาสุ่ม การลดรูปแลตทิซ และการโจมตีแบบผสมผสาน}
	\end{quote}
	
	ดังนั้น เอกสารฉบับนี้สามารถใช้เป็นต้นฉบับอ้างอิงเพื่อ
	\begin{enumerate}[label=\arabic*.]
		\item ทำรายงานภาษาไทยฉบับสมบูรณ์,
		\item ดัดแปลงเป็นบทพูดหรือบันทึกประกอบการประชุม,
		\item และใช้สร้าง \eng{presentation} เชิงเทคนิคได้อย่างมีโครงสร้างและครบถ้วน
	\end{enumerate}
	
	% =========================================================
	\section*{ภาคผนวก: สรุปคำสำคัญที่ควรใช้ในงานนำเสนอ}
	\addcontentsline{toc}{section}{ภาคผนวก: สรุปคำสำคัญที่ควรใช้ในงานนำเสนอ}
	
	\begin{longtable}{p{0.28\textwidth} p{0.62\textwidth}}
		\toprule
		คำสำคัญ & คำอธิบายแบบสั้นสำหรับใส่ในสไลด์ \\
		\midrule
		\eng{Masked-Salting} & การเพิ่มเมทริกซ์เบาบาง $M$ ลงใน $A$ เพื่อสร้างความไม่แน่นอนเชิงโครงสร้าง \\
		\eng{Indistinguishability} & ผู้โจมตีไม่สามารถแยกแยะระบบใหม่ออกจากระบบเดิมได้ง่ายในเชิงสถิติ \\
		\eng{Exhaustive Search} & การเดาแบบครอบจักรวาล ใช้เป็นฐานวัดความซับซ้อนขั้นต่ำ \\
		\eng{BKZ Reduction} & เทคนิคการลดรูปแลตทิซที่ใช้วิเคราะห์ความได้เปรียบเชิงอัลกอริทึมของผู้โจมตี \\
		\eng{Hybrid Attack} & การใช้ \eng{BKZ} ช่วยคัดกรองก่อน แล้วจึง \eng{Brute-force} เฉพาะส่วนที่เหลือ \\
		\eng{Weak Lattice} & โครงสร้างแลตทิซที่อาจทำให้ผู้โจมตีลดต้นทุนของการถอดรหัสได้ผิดปกติ \\
		\eng{Exponential Complexity} & ความซับซ้อนที่เติบโตเร็วมากเมื่อมิติของปัญหาเพิ่มขึ้น \\
		\eng{Computational Infeasibility} & ต้นทุนการโจมตีสูงจนไม่สามารถทำได้จริงในทางปฏิบัติ \\
		\bottomrule
	\end{longtable}
	
\end{document}
